% Recuperacion de 09c_entropia_sectorial.tex, 11_confluencia.tex y
% 85_agujero_negro_doble_hoja.tex; originales en fuentes_conservadas/
% horizontes_informacion. Procedencia: desarrollo/HORIZONTES_INFORMACION_FUENTES.md.
% Las referencias a antecedentes se acuerdan con el editor de VII.
% No incorpora una corriente material ni evalua la amplitud torsional.

\section{Horizons, purification et bilans d'information}
\label{vii:sec:horizontes-informacion}

L'état enrichi conserve les parcours et leurs registres pendant l'évolution ; une lecture de sous-système détermine quelle partie de cette information reste accessible. La construction de cette section utilise les canaux angulaires de la section~\ref{v:ant:sec:ii-angulos}, l'incidence et l'opérateur d'aire de la section~\ref{v:ant:sec:area}, ainsi que les unités d'action, d'aire et de température des sections~\ref{v:ant:sec:gravity} et~\ref{v:ant:sec:ii-kb-aut-desarrollo}. Ce sont des représentations antérieures de la même construction APP--TRIT--TPK. Chacune transmet ici son domaine, son orientation, sa normalisation et ses conditions de réalisation.

On distingue trois espaces : le registre binaire de feuillet, les facteurs
tritiques de l'incidence sectorielle et la réalisation géométrique des
extérieurs d'un horizon. La purification, le transport des observables
et la conversion dimensionnelle précisent leurs relations. Cette
distinction conserve la structure qui accompagne chaque valeur entropique.

\subsection{Registre de feuillet et purification bipartite}
\label{vii:subsec:hojas-purificacion}

Écrivons \(a=A_{\rm rad}\) et \(y=C^*_{\rm rad}\) pour les
coordonnées angulaires déjà construites. Les canaux et leurs probabilités sont
\begin{equation}
 \mathfrak q_\pm=e^{-a\mp y},\qquad
 p_\pm=\frac{\mathfrak q_\pm}{\mathfrak q_++\mathfrak q_-}
       =\frac{e^{\mp y}}{2\cosh y}.
 \label{vii:eq:probabilidades-hoja}
\end{equation}
Les signes indexent ces deux canaux. Dans
\(\mathcal H_{\rm h}=\operatorname{span}\{|+\rangle,|-\rangle\}\),
nous définissons
\[
 \rho_{\rm h}(y)=p_+|+\rangle\langle+|+p_-|-\rangle\langle-|.
\]
L'espace auxiliaire \(\mathcal H_{\rm ar}=\operatorname{span}
\{|\mathrm{adv}\rangle,|\mathrm{ret}\rangle\}\) porte deux étiquettes orthonormées. Une phase \(\chi_0\) étant fixée, l'état conjoint est
\begin{equation}
 |\Psi_{\rm ar}\rangle
 =\sqrt{p_+}|+\rangle\otimes|\mathrm{adv}\rangle
  +e^{i\chi_0}\sqrt{p_-}|-\rangle\otimes|\mathrm{ret}\rangle.
 \label{vii:eq:purificacion-hoja}
\end{equation}

\begin{proposition}[Information accessible dans le registre de feuillet]
\label{vii:prop:purificacion-hoja}
Le vecteur~\eqref{vii:eq:purificacion-hoja} est normalisé et sa trace
partielle sur \(\mathcal H_{\rm ar}\) est \(\rho_{\rm h}(y)\).
Son entropie d'intrication, exprimée en nats, est
\begin{equation}
 s_{\rm h}(y)=\log(2\cosh y)-y\tanh y.
 \label{vii:eq:entropia-hoja}
\end{equation}
L'échange des étiquettes \(+\) et \(-\) conjugue
\(\rho_{\rm h}(y)\) à \(\rho_{\rm h}(-y)\) et conserve
\(s_{\rm h}\).
\end{proposition}
\begin{proof}
Le carré de la norme est \(p_++p_-=1\). L'orthogonalité des étiquettes auxiliaires élimine les termes croisés de la trace partielle. Les carrés des coefficients de Schmidt sont précisément \(p_\pm\). Par~\eqref{vii:eq:probabilidades-hoja}, \(p_+-p_-=-\tanh y\) ; substituer \(\log p_\pm=\mp y-\log(2\cosh y)\) dans \(-\sum_\pm p_\pm\log p_\pm\) démontre la formule. L'échange des signes permute les deux valeurs propres et conserve l'entropie.
\end{proof}

Le facteur commun \(e^{-a}\) se simplifie lors de la normalisation du registre binaire.
La phase relative \(\chi_0\) demeure dans l'état conjoint bien que cette
réduction diagonale ne la publie pas. La réversibilité logique se rapporte
au transport du registre complet, avec ses étiquettes et ses corrélations ;
le retour de phase conserve l'avancée de la mémoire décrite dans la
section~\ref{vii:sec:retorno-memoria-gravedad}. Une mise en œuvre thermique
a en outre sa propre loi d'énergie et de dissipation. En particulier,
\(s_{\rm h}\) mesure ce sous-système binaire et sa bipartition précise.

\subsection{Purification des secteurs d'incidence}
\label{vii:subsec:purificacion-sectorial}

L'incidence de bord construite précédemment distingue deux ensembles
de dix-huit liens, d'étiquettes sectorielles \(m=90,120\). Chaque lien
porte \(\mathbb C^3\), de base \(|0\rangle,|1\rangle,|2\rangle\)
et d'opérateur d'occupation \(\operatorname{diag}(0,1,2)\). Soient
\(\mathsf N_m\) leurs sommes sectorielles, et
\begin{equation}
 \mathsf N=\mathsf N_{90}+\mathsf N_{120},\qquad
 \mathsf D_\partial=90\mathsf N_{90}+120\mathsf N_{120}.
 \label{vii:eq:ocupaciones-sectoriales}
\end{equation}
Dans la branche positive, \(\gamma_{\rm BI}\) désigne le paramètre de Barbero--Immirzi déjà obtenu, et \(a_0>0\) son facteur angulaire d'aire. Cette notation distingue le paramètre de la connexion et de ses holonomies. L'aire s'exprime comme
\begin{equation}
 \widehat{\mathcal A}_m=\ell_P^2\gamma_{\rm BI}a_0\mathsf N_m,
 \qquad
 \widehat{\mathcal A}=\widehat{\mathcal A}_{90}
                         +\widehat{\mathcal A}_{120}.
 \label{vii:eq:area-sectores}
\end{equation}

Pour le paramètre fini \(\Delta=\Delta_{\rm mod}\in\mathbb R\) de l'état réduit, nous posons
\begin{align}
 x_m&=e^{-m\Delta},& z_m&=1+x_m+x_m^2,&
 Z(\Delta)&=z_{90}^{18}z_{120}^{18},\\
 \rho_\partial(\Delta)&=Z(\Delta)^{-1}e^{-\Delta\mathsf D_\partial}.
 \label{vii:eq:estado-sectorial}
\end{align}
La trace dans cette section est la trace matricielle ordinaire et les états ont une trace égale à un. Le paramètre \(\Delta\) conserve la sélection et la normalisation de l'état antécédent.

\begin{proposition}[Purification tensorielle de la frontière]
\label{vii:prop:purificacion-sectorial}
Pour un poids réel fini \(w\), soit
\begin{equation}
 |\operatorname{TFD}(w)\rangle
 =\frac{|0\rangle|0\rangle+e^{-w/2}|1\rangle|1\rangle
                   +e^{-w}|2\rangle|2\rangle}
        {\sqrt{1+e^{-w}+e^{-2w}}}.
 \label{vii:eq:tfd-local}
\end{equation}
Le produit
\begin{equation}
 |\Psi_\Delta\rangle
 =|\operatorname{TFD}(90\Delta)\rangle^{\otimes18}
  \otimes|\operatorname{TFD}(120\Delta)\rangle^{\otimes18}
 \label{vii:eq:tfd-sectorial}
\end{equation}
est une purification normalisée de \(\rho_\partial(\Delta)\).
\end{proposition}
\begin{proof}
Dans chaque lien, la somme des carrés des trois coefficients vaut un.
La trace partielle élimine les termes d'indices distincts et restitue
\(\operatorname{diag}(1,e^{-w},e^{-2w})/(1+e^{-w}+e^{-2w})\).
La trace partielle commute avec le produit tensoriel. Multiplier les
trente-six réductions donne l'exponentielle et la fonction de partition
de~\eqref{vii:eq:estado-sectorial}.
\end{proof}

\subsection{Équation informationnelle et fluctuations surfaciques}
\label{vii:subsec:ecuacion-informacional}

Les moyennes et les variances d'occupation par lien sont
\begin{equation}
 f_m=\frac{x_m+2x_m^2}{z_m},\qquad
 v_m=\frac{x_m+4x_m^2}{z_m}-f_m^2.
 \label{vii:eq:momentos-sectoriales}
\end{equation}
Nous écrivons \(s(\rho)=-\operatorname{Tr}(\rho\log\rho)\), avec \(0\log0=0\). Dans la purification précédente, \(s(\rho_\partial)\) est son entropie d'intrication.

\begin{proposition}[Aire moyenne, entropie et réponse sectorielle]
\label{vii:prop:respuesta-sectorial}
Pour \(\Delta\) réel et fini, on a
\begin{align}
 \langle\mathsf N\rangle_\Delta&=18(f_{90}+f_{120}),\\
 \langle\mathsf D_\partial\rangle_\Delta&=18(90f_{90}+120f_{120}),\\
 \langle\widehat{\mathcal A}\rangle_\Delta
 &=18\ell_P^2\gamma_{\rm BI}a_0(f_{90}+f_{120}),\\
 s(\Delta)&=\log Z(\Delta)+\Delta\langle\mathsf D_\partial\rangle_\Delta,
 \label{vii:eq:entropia-sectorial}\\
 \frac{\mathrm d\langle\mathsf D_\partial\rangle_\Delta}{\mathrm d\Delta}
 &=-\operatorname{Var}_\Delta(\mathsf D_\partial)
  =-18(90^2v_{90}+120^2v_{120}),\\
 \frac{\mathrm ds}{\mathrm d\Delta}
 &=-\Delta\operatorname{Var}_\Delta(\mathsf D_\partial),\\
 \frac{\mathrm d\langle\widehat{\mathcal A}\rangle_\Delta}{\mathrm d\Delta}
 &=-18\ell_P^2\gamma_{\rm BI}a_0(90v_{90}+120v_{120}).
 \label{vii:eq:respuestas-area-entropia}
\end{align}
\end{proposition}
\begin{proof}
Les poids locaux sont \((1,x_m,x_m^2)/z_m\), dont les deux premiers moments donnent \(f_m\) et \(v_m\). L'additivité des opérateurs d'occupation et la factorisation de l'état donnent les trois premières formules. Prendre la trace de \(-\rho_\partial\log\rho_\partial
=\rho_\partial(\log Z+\Delta\mathsf D_\partial)\) donne l'entropie. La différentiation des sommes finies produit \(\partial_\Delta\log Z=-\langle\mathsf D_\partial\rangle_\Delta\) et \(\partial_\Delta^2\log Z
=\operatorname{Var}_\Delta(\mathsf D_\partial)\). Les facteurs indépendants additionnent leurs variances. Dans la dérivée de \eqref{vii:eq:entropia-sectorial}, les deux termes de premier moment s'annulent. Enfin, \(\partial_\Delta f_m=-mv_m\) donne la réponse de l'aire.
\end{proof}

La perte d'uniformité se quantifie en dimension \(d_\partial=3^{36}\) :
\begin{equation}
 I_\partial(\Delta)
 =36\log3-s(\Delta)
 =\mathcal D\bigl(\rho_\partial(\Delta)\Vert I/d_\partial\bigr)\ge0.
 \label{vii:eq:deficit-informacional}
\end{equation}
Ici \(\mathcal D\) désigne l'entropie relative, définie et justifiée ci-dessous. La substitution de \(\log(I/d_\partial)=-36\log3\,I\) démontre l'identité. La convexité de \(t\log t\) donne la positivité ; l'égalité correspond à la distribution uniforme, qui, dans cette collection, se produit exactement en \(\Delta=0\). Le déficit mesure une information adimensionnelle ; l'échelle \(\gamma_{\rm BI}a_0\ell_P^2\) mesure une aire par occupation.

\begin{theorem}[Inversion des occupations et aire complémentaire]
\label{vii:thm:inversion-sectorial}
Soit \(\mathcal R_\partial\) l'involution qui échange
\(|0\rangle\) et \(|2\rangle\) et fixe \(|1\rangle\) dans chaque
facteur. Alors
\begin{align}
 \mathcal R_\partial\mathsf N\mathcal R_\partial^*&=72I-\mathsf N,\\
 \mathcal R_\partial\mathsf D_\partial\mathcal R_\partial^*
 &=7560I-\mathsf D_\partial,\\
 \rho_\partial(-\Delta)&=\mathcal R_\partial\rho_\partial(\Delta)
                          \mathcal R_\partial^*,\\
 s(-\Delta)&=s(\Delta),\\
 \langle\widehat{\mathcal A}\rangle_{-\Delta}
 +\langle\widehat{\mathcal A}\rangle_\Delta
 &=72\gamma_{\rm BI}a_0\ell_P^2,\\
 \mathcal D\bigl(\rho_\partial(\Delta)\Vert\rho_\partial(-\Delta)\bigr)
 &=\Delta\bigl(7560-2\langle\mathsf D_\partial\rangle_\Delta\bigr).
 \label{vii:eq:inversion-sectorial}
\end{align}
\end{theorem}
\begin{proof}
Chaque occupation se transforme par \(n\mapsto2-n\). Dans chaque secteur,
\(\mathsf N_m\mapsto36I-\mathsf N_m\), de sorte que le terme
constant du lecteur pondéré est \(36(90+120)=7560\).
De plus, \(z_m(-\Delta)=e^{2m\Delta}z_m(\Delta)\) et
\(Z(-\Delta)=e^{7560\Delta}Z(\Delta)\). La substitution dans l'
état normalisé prouve la conjugaison. Son spectre conserve l'entropie ;
la transformation de \(\mathsf N\) donne l'aire complémentaire. Soustraire
les logarithmes des deux états et prendre l'espérance dans
\(\rho_\partial(\Delta)\) donne la dernière égalité.
\end{proof}

En \(\Delta=0\), \(s=36\log3\). Lorsque \(\Delta\to+\infty\),
l'occupation nulle concentre l'état et l'entropie tend vers zéro ;
l'inversion envoie cette limite sur l'occupation maximale. Pour \(\Delta>0\),
l'entropie décroît strictement parce que le lecteur pondéré possède plus
d'une valeur propre et que l'état est de rang plein. La configuration
complète, l'aire et l'entropie conservent ainsi des fonctions différenciées.

\subsection{Loi modulaire et variations finies d'aire}
\label{vii:subsec:ley-modular-finita}

Fixons \(\rho_0=\rho_\partial(\Delta_0)\), avec \(\Delta_0\) réel fini, et maintenons \(\Delta_0,\gamma_{\rm BI},a_0,\ell_P\) constants. Les aires normalisées \(\mathcal A_m^{\rm n}=\gamma_{\rm BI}a_0\mathsf N_m\) expriment l'opérateur modulaire comme
\begin{equation}
 \mathsf K_0=-\log\rho_0
 =c_0I+\sum_{m=90,120}\beta_m\mathcal A_m^{\rm n},
 \quad \beta_m=\frac{m\Delta_0}{\gamma_{\rm BI}a_0},
 \quad c_0=\log Z(\Delta_0).
 \label{vii:eq:operador-modular-areal}
\end{equation}
On a \(\beta_{120}=\tfrac43\beta_{90}\), y compris lorsque les deux sont nuls. Le quotient est \(4/3\) lorsque \(\Delta_0\ne0\). \(\mathsf K_0\) est le logarithme de cet état de référence ; son type le distingue du transport temporel et de la connexion gravitationnelle.

\begin{theorem}[Première loi modulaire et correction finie]
\label{vii:thm:ley-modular-finita}
Pour tout état normalisé \(\sigma\) du même espace, soit
\(\Delta_\sigma\langle B\rangle
=\operatorname{Tr}[(\sigma-\rho_0)B]\). Alors
\begin{equation}
 s(\sigma)-s(\rho_0)
 =\sum_{m=90,120}\beta_m\Delta_\sigma\langle\mathcal A_m^{\rm n}\rangle
  -\mathcal D(\sigma\Vert\rho_0),
 \label{vii:eq:ley-modular-finita}
\end{equation}
où
\(\mathcal D(\sigma\Vert\rho_0)
=\operatorname{Tr}[\sigma(\log\sigma-\log\rho_0)]\)
est finie et non négative. Si \(\sigma(\varepsilon)\) est différentiable,
normalisée et \(\sigma(0)=\rho_0\), on obtient
\begin{equation}
 \left.\frac{\mathrm ds(\sigma(\varepsilon))}{\mathrm d\varepsilon}\right|_0
 =\sum_{m=90,120}\beta_m
   \left.\frac{\mathrm d\langle\mathcal A_m^{\rm n}\rangle_{
                 \sigma(\varepsilon)}}{\mathrm d\varepsilon}\right|_0.
 \label{vii:eq:primera-ley-modular}
\end{equation}
\end{theorem}
\begin{proof}
Par définition,
\(\mathcal D(\sigma\Vert\rho_0)=-s(\sigma)
+\operatorname{Tr}(\sigma\mathsf K_0)\), tandis que
\(s(\rho_0)=\operatorname{Tr}(\rho_0\mathsf K_0)\).
La soustraction et~\eqref{vii:eq:operador-modular-areal} prouvent l'identité
finie, car la normalisation annule le terme \(c_0I\).

Pour la positivité, diagonalisons \(\sigma\) avec les valeurs propres \(p_i\) et les vecteurs \(u_i\), et \(\rho_0\) avec les valeurs propres \(q_j>0\) et les vecteurs \(v_j\). Si \(r_i=\langle u_i,\rho_0u_i\rangle>0\), la concavité du logarithme donne
\[
 \sum_j|\langle u_i,v_j\rangle|^2\log q_j\le\log r_i.
\]
Par conséquent,
\(\mathcal D(\sigma\Vert\rho_0)
\ge\sum_i p_i\log(p_i/r_i)\ge0\).
La dernière inégalité est la convexité de \(t\log t\) avec les poids
\(r_i\), dont la somme vaut un. La fidélité de \(\rho_0\) maintient
ses logarithmes finis ; les valeurs propres nulles de \(\sigma\)
s'interprètent au moyen de \(0\log0=0\).

Pour la formule infinitésimale, au voisinage de l'état fidèle \(\rho_0\),
la dérivée de la trace de \(-\sigma\log\sigma\) est
\(-\operatorname{Tr}[\dot\sigma(\log\rho_0+I)]\).
Cette identité s'obtient à partir de la dérivée spectrale de la trace.
La condition \(\operatorname{Tr}\dot\sigma=0\) laisse
\(\operatorname{Tr}(\dot\sigma\mathsf K_0)\), et substituer
\eqref{vii:eq:operador-modular-areal} conclut la preuve.
\end{proof}

Pour les aires physiques, les coefficients sont \(\beta_m/\ell_P^2\). L'application du transducteur commun \(k_B\) donne \(S=k_Bs\) et conserve intégralement le terme \(k_B\mathcal D\). La dérivée dans \eqref{vii:eq:respuestas-area-entropia} fait varier \(\Delta\) ; la loi \eqref{vii:eq:primera-ley-modular} maintient \(\Delta_0\) fixe et fait varier l'état autour de la référence. Ce sont des opérations distinctes.

\subsection{Action, énergie de décision et unité surfacique}
\label{vii:subsec:conversion-informacion-area}

Nous fixons une même section orientée d'action et ses sections gravitationnelle
et thermique, sur la branche positive. La composition
de~\ref{g26:sec:rigidez-radial-es} identifie
\(\ell_P=L_\alpha=L\) dans cette réalisation. Dans la section de retour,
l'énergie caractéristique est \(E_L=\hbar_{\rm ret}c/L\). On a
\begin{equation}
 \begin{aligned}
 \ell_P^2&=\frac{\hbar G}{c^3},& t_P&=\frac{\ell_P}{c},\\
 T_{108}&=2\pi t_P,& h&=2\pi\hbar,\\
 E_P&=\frac{\hbar}{t_P}=k_B\Theta_{\rm clk}\log3.&&
 \end{aligned}
 \label{vii:eq:unidad-accion-informacion}
\end{equation}
L'égalité thermique conserve le couple et la normalisation précédemment
construits ; lors du changement de leur base dimensionnelle, on transporte conjointement
\(k_B\) et \(\Theta_{\rm clk}\). L'incrément élémentaire d'aire est
\(\delta\mathcal A=\gamma_{\rm BI}a_0\ell_P^2\).

\begin{proposition}[Coefficient de conversion entre décision et aire]
\label{vii:prop:tension-informacion-area}
Le quotient de l'énergie de décision par l'incrément surfacique satisfait
\begin{equation}
 \begin{aligned}
 \mathcal T_{I/A}:=\frac{E_P}{\delta\mathcal A}
 &=\frac{k_B\Theta_{\rm clk}\log3}{\gamma_{\rm BI}a_0\ell_P^2}\\
 &=\frac{c^3}{\gamma_{\rm BI}a_0G t_P}
  =\frac{c^4}{\gamma_{\rm BI}a_0G\ell_P}.
 \end{aligned}
 \label{vii:eq:tension-informacion-area}
\end{equation}
Sa dimension est celle d'une énergie par aire.
\end{proposition}
\begin{proof}
La substitution de \(E_P=\hbar/t_P\) et de
\(\ell_P^2=\hbar G/c^3\) annule la section d'action commune.
L'identité \(\ell_P=ct_P\) donne la dernière expression. La positivité
des facteurs garantit que tous les quotients sont définis.
\end{proof}

La capacité \(81\) de la coupe cubique et la bipartition des trente-six liaisons sectorielles correspondent à des domaines différents. Dans le graphe cubique \(\{0,\ldots,N-1\}^3\), les \(N^2\) droites parallèles entre deux faces opposées sont disjointes par les liaisons : toute coupe intercepte les \(N^2\), et une couche transverse atteint cette borne. Pour \(N=9\), l'état produit uniforme des trits de cette couche possède \(s_{\rm corte}=81\log3\). L'attribution d'une décision du même cycle à chaque liaison exprime
\begin{equation}
 S_{\rm corte}^{\rm fis}=81k_B\log3,\qquad
 E_{\rm corte}^{(1)}=81E_P
                  =\Theta_{\rm clk}S_{\rm corte}^{\rm fis}.
 \label{vii:eq:corte-energia-informacion}
\end{equation}
La preuve combine la coupe minimale, l'additivité de l'entropie de l'état produit et~\eqref{vii:eq:unidad-accion-informacion}. Dans l'état sectoriel, l'équation finie pertinente reste \eqref{vii:eq:ley-modular-finita}, avec son entropie relative.

\subsection{Normalisation cosmologique de l'horizon}
\label{vii:subsec:balance-horizonte}

Pour un rayon aréolaire \(R>0\), la réalisation cosmologique considérée utilise la triade
\begin{equation}
 E_\Lambda(R)=\frac{c^4R}{2G},\qquad
 S_{\rm H}(R)=\frac{k_B\pi R^2}{\ell_P^2},\qquad
 T_{\rm cos}(R)=\frac{\hbar c}{2\pi k_BR}.
 \label{vii:eq:horizonte-cosmologico}
\end{equation}
Le facteur \(2\pi\) appartient à cette normalisation de température.
L'aire \(\mathcal A_{\rm H}=4\pi R^2\) exprime la même entropie
sous la forme \(S_{\rm H}/k_B=\mathcal A_{\rm H}/(4\ell_P^2)\).

Pour conserver la différence entre échelle construite et rayon variable, nous écrivons ici \(R=R_h=\zeta_hL_\alpha\), avec \(\zeta_h>0\). Dans la section de retour, substituer le même \(G\) et la même aire de référence donne
\[
 E_\Lambda(R_h)=\frac{\zeta_h}{2}E_L,\qquad
 \frac{S_{\rm H}(R_h)}{k_B}=\pi\zeta_h^2,\qquad
 k_BT_{\rm cos}(R_h)=\frac{E_L}{2\pi\zeta_h}.
\]
Chaque égalité procède directement de
\eqref{vii:eq:horizonte-cosmologico}. Le facteur \(1/2\) appartient
à cette énergie d'horizon. La lecture réduite d'une masse
\(m_{P,\rm ret}y\) a pour rayon \(r_g=L_\alpha y\) ; son rayon de
Schwarzschild est \(2L_\alpha y\). La condition géométrique d'un
horizon cosmologique maintient son domaine propre.

\begin{theorem}[Cellule gravitationnelle de demi-action]
\label{vii:thm:celda-media-accion}
Dans~\eqref{vii:eq:horizonte-cosmologico},
\begin{align}
 T_{\rm cos}(R)S_{\rm H}(R)&=E_\Lambda(R),\\
 T_{\rm cos}(\ell_P)S_{\rm H}(\ell_P)
 &=\frac{E_P}{2}=\frac{k_B\Theta_{\rm clk}\log3}{2},\\
 E_\Lambda(\ell_P)t_P&=\frac{\hbar}{2},&
 E_\Lambda(\ell_P)T_{108}&=\frac h2.
 \label{vii:eq:media-accion-horizonte}
\end{align}
En maintenant fixes les constantes de la coordinatisation et en faisant varier \(R\),
\begin{equation}
 T_{\rm cos}\,\mathrm dS_{\rm H}=2\,\mathrm dE_\Lambda,
 \qquad S_{\rm H}\,\mathrm dT_{\rm cos}=-\mathrm dE_\Lambda.
 \label{vii:eq:diferencial-horizonte}
\end{equation}
\end{theorem}
\begin{proof}
Multiplier la température par l'entropie donne
\(\hbar cR/(2\ell_P^2)=c^4R/(2G)\).
En \(R=\ell_P\), le résultat est
\(\hbar/(2t_P)=E_P/2\). Les deux durées de
\eqref{vii:eq:unidad-accion-informacion} donnent les relations d'action.
Pour les différentielles,
\(\mathrm dE_\Lambda=c^4\mathrm dR/(2G)\),
\(\mathrm dS_{\rm H}=2\pi k_BR\mathrm dR/\ell_P^2\) et
\(\mathrm dT_{\rm cos}=-T_{\rm cos}\mathrm dR/R\).
La substitution prouve les deux identités.
\end{proof}

Le produit \(E_\Lambda=T_{\rm cos}S_{\rm H}\) exprime une identité d'état. Ses variations radiales conservent les deux termes de~\eqref{vii:eq:diferencial-horizonte} ; une interprétation thermodynamique complète spécifie en outre le générateur temporel, le travail et la réalisation matérielle correspondants. Dans cette coordinatisation, l'action et la gravitation fixent l'unité \(\ell_P^2\), tandis que Boltzmann convertit l'information en entropie dimensionnelle. Le transport des sections avec \(G\hbar\) fixé conserve cette unité lorsque la section \(c\) reste fixe, et conserve \(S_{\rm H}/k_B\) à rayon fixé.

Dans la cellule \(R=\ell_P\), \(S_{\rm H}/k_B=\pi\), de sorte que
la multiplicité effective est \(\Omega_{\rm eff}=e^\pi\).
Pour l'involution hyperbolique tritique \(H_{\rm trit}\), de spectre
\(\{1,-1\}\), la diagonalisation donne
\(\operatorname{spec}(e^{\pi H_{\rm trit}})=\{e^\pi,e^{-\pi}\}\).
Cette relation spectrale spécifie la valeur expansive ; une identification
microscopique des deux espaces conserve en outre son application de
réalisation et l'état.

% Ampliación aditiva: acción, área, memoria sectorial y termometría.
% Propietarios: 70_horizontes_e_informacion.tex; sección térmica antecedente;
% 71_energia_libre_holografica_cerrada.tex del corpus integral.
% Se mantienen las condiciones de la realización cosmológica declarada.

\subsection{Conversion thermique entre le calendrier et l'horizon}
\label{viii:subsec:conversion-termica-calendario-horizonte}

L'action, l'incrément d'aire et la mémoire angulaire admettent une composition énergétique explicite. On conserve une même orientation d'action, la branche positive de Barbero--Immirzi et la réalisation de \eqref{vii:eq:horizonte-cosmologico}. Nous abrégeons
\[
 L=\ell_P=L_\alpha,\qquad
 E_P=\frac{h}{108t_0}=\frac{\hbar c}{L},\qquad
 q_{\mathcal A}=\delta\mathcal A=\gamma_{\rm BI}a_0L^2.
\]
Le symbole \(q_{\mathcal A}\) désigne ici un incrément d'aire.
Le facteur angulaire \(a_0\), la fonctionnelle \(\gamma_{\rm BI}\) et
l'action sont les grandeurs obtenues dans les sections précédentes.
Leur domaine conserve l'incidence hexadique--octadique, l'orientation et
le calendrier du même état enrichi APP--TRIT--TPK.

\begin{proposition}[Composition des échelles surfacique et thermique]
\label{viii:prop:composicion-area-termica}
Avec la tension de~\eqref{vii:eq:tension-informacion-area}, on a
\begin{equation}
 \frac{\mathcal T_{I/A}q_{\mathcal A}}{\log3}
 =\frac{E_P}{\log3}
 =\frac{h}{108t_0\log3}
 =k_B\Theta_{\rm clk}.
 \label{viii:eq:composicion-area-termica}
\end{equation}
Pour \(R=\zeta L\), avec \(\zeta>0\), soit
\(\tau_R=k_BT_{\rm cos}(R)\). Alors
\begin{equation}
 \tau_R=\frac{E_P}{2\pi\zeta},\qquad
 \frac{T_{\rm cos}(R)}{\Theta_{\rm clk}}
 =\frac{\log3}{2\pi\zeta}.
 \label{viii:eq:conversion-dos-temperaturas}
\end{equation}
En particulier, le quotient dans \(R=L\) est \(\log3/(2\pi)\).
\end{proposition}
\begin{proof}
La définition \(\mathcal T_{I/A}=E_P/q_{\mathcal A}\) et la
normalisation tritique de~\eqref{vii:eq:unidad-accion-informacion}
donnent~\eqref{viii:eq:composicion-area-termica}. L'expression
gravitationnelle de la même tension conduit au même résultat :
\[
 \frac{c^4}{\gamma_{\rm BI}a_0GL}
        \gamma_{\rm BI}a_0L^2
 =\frac{c^4L}{G}=\frac{\hbar c}{L}=E_P,
\]
où l'on a utilisé \(G\hbar=c^3L^2\).
D'après~\eqref{vii:eq:horizonte-cosmologico},
\(\tau_R=\hbar c/(2\pi R)=E_P/(2\pi\zeta)\).
Diviser cette égalité par \(k_B\Theta_{\rm clk}=E_P/\log3\)
prouve le quotient des températures.
\end{proof}

Dans la même section d'action, la composition gravitationnelle s'exprime par
\begin{equation}
 Gk_B\Theta_{\rm clk}\log3=GE_P=c^4L,
 \qquad G\tau_R=\frac{c^4L}{2\pi\zeta}.
 \label{viii:eq:composicion-termica-gravitatoria}
\end{equation}
Les deux identités découlent de \(G\hbar=c^3L^2\) et des échelles précédentes.

La tension convertit l'incrément surfacique en énergie de décision.
La normalisation tritique et la périodicité cosmologique déterminent deux
sections thermiques de cette même énergie, reliées par
\eqref{viii:eq:conversion-dos-temperaturas}. La conversion conserve
les deux normalisations et la fonction de Barbero--Immirzi dans la résolution
géométrique.

\subsection{Réponse énergétique de la mémoire sectorielle}
\label{viii:subsec:respuesta-energetica-sectorial}

Fixons \(R>0\) et l'état fidèle
\(\rho_0=\rho_\partial(\Delta_0)\), avec \(\Delta_0>0\),
de~\eqref{vii:eq:estado-sectorial}. Son opérateur modulaire est
\(\mathsf K_0=(\log Z_0)I+\Delta_0\mathsf D_\partial\).
La réalisation d'horizon associe à chaque état de densité
\(\sigma\) les grandeurs de réponse
\begin{equation}
\begin{aligned}
 \mathcal E_R(\sigma)
 &=E_\Lambda(R)+\tau_R
   \operatorname{Tr}[(\sigma-\rho_0)\mathsf K_0],\\
 \mathcal S_R(\sigma)
 &=S_{\rm H}(R)+k_B[s(\sigma)-s(\rho_0)].
\end{aligned}
 \label{viii:eq:respuesta-horizonte-definida}
\end{equation}
Il s'agit de la réponse de l'état modulaire dans la réalisation géométrique déjà spécifiée. Le rayon et l'action fixent son échelle énergétique \(\tau_R\) ; les occupations conservent la provenance de chaque canal.

\Needspace{28\baselineskip}
\begin{theorem}[Hamiltonien de réponse et bilan informationnel]
\label{viii:thm:hamiltoniano-respuesta-termica}
L'opérateur des différences énergétiques correspondant à \eqref{viii:eq:respuesta-horizonte-definida} peut être choisi comme
\begin{equation}
 \mathsf H_R=\tau_R\Delta_0\mathsf D_\partial.
 \label{viii:eq:hamiltoniano-respuesta}
\end{equation}
Son état de Gibbs à l'échelle énergétique \(\tau_R\) est \(\rho_0\).
Pour tout état de densité \(\sigma\) dans l'espace sectoriel,
\begin{equation}
 \mathcal E_R(\sigma)-T_{\rm cos}(R)\mathcal S_R(\sigma)
 =\tau_R\mathcal D(\sigma\Vert\rho_0)\geq0.
 \label{viii:eq:balance-libre-termico}
\end{equation}
Dans une variation différentiable des états autour de \(\rho_0\),
avec \(R\) et \(\mathsf H_R\) fixés, on obtient
\begin{equation}
 \delta\mathcal E_R
 =\tau_R\,\delta s
 =T_{\rm cos}(R)\,\delta\mathcal S_R.
 \label{viii:eq:primera-ley-energetica}
\end{equation}
\end{theorem}
\begin{proof}
Les traces de \(\sigma\) et de \(\rho_0\) valent un, de sorte que la différence élimine le terme scalaire \(\log Z_0\). Il en résulte \(\mathcal E_R(\sigma)-E_\Lambda(R)
=\operatorname{Tr}[(\sigma-\rho_0)\mathsf H_R]\). En outre,
\[
 \frac{e^{-\mathsf H_R/\tau_R}}
 {\operatorname{Tr}e^{-\mathsf H_R/\tau_R}}
 =\frac{e^{-\Delta_0\mathsf D_\partial}}{Z_0}=\rho_0.
\]
L'égalité \(E_\Lambda=T_{\rm cos}S_{\rm H}\) annule la
référence géométrique dans le bilan. L'identité finie
\eqref{vii:eq:ley-modular-finita}, écrite sous la forme
\(s(\sigma)-s(\rho_0)
=\operatorname{Tr}[(\sigma-\rho_0)\mathsf K_0]
-\mathcal D(\sigma\Vert\rho_0)\), démontre
\eqref{viii:eq:balance-libre-termico}. Son terme tangent est la
première loi modulaire déjà prouvée et produit
\eqref{viii:eq:primera-ley-energetica}.
\end{proof}

Pour un incrément unitaire d'occupation, les sauts énergétiques sont
\(\epsilon_{90}=90\tau_R\Delta_0\) et
\(\epsilon_{120}=120\tau_R\Delta_0\). Dans les deux secteurs,
\begin{equation}
 \frac{\epsilon_m}{m\Delta_0}=\tau_R,
 \qquad m\in\{90,120\}.
 \label{viii:eq:pendiente-termica-comun}
\end{equation}
Ainsi, le quotient \(4/3\) entre poids surfaciques enregistre les différents sauts du même hamiltonien. Le quotient entre un saut énergétique et la variation correspondante de \(-\log p\) conserve une échelle commune. La normalisation de l'état n'élimine qu'une constante scalaire.

\begin{corollary}[Susceptibilité énergétique à hamiltonien fixé]
\label{viii:cor:respuesta-escala-termica}
Pour \(\rho_\tau=e^{-\mathsf H_R/\tau}/
\operatorname{Tr}e^{-\mathsf H_R/\tau}\), avec \(\tau>0\),
\begin{equation}
 \frac{\mathrm d\langle\mathsf H_R\rangle_\tau}{\mathrm d\tau}
 =\frac{\operatorname{Var}_{\rho_\tau}(\mathsf H_R)}{\tau^2}>0.
 \label{viii:eq:susceptibilidad-termica}
\end{equation}
En \(\tau=\tau_R\), le résultat est
\(\Delta_0^2\operatorname{Var}_{\rho_0}(\mathsf D_\partial)\).
\end{corollary}
\begin{proof}
La dérivation de la somme finie d'exponentielles donne la variance divisée
par \(\tau^2\). L'état est fidèle et \(\mathsf H_R\) possède plus d'
une valeur propre, de sorte que la variance est positive. La substitution de
\eqref{viii:eq:hamiltoniano-respuesta} donne l'évaluation indiquée.
\end{proof}

Cette dérivée fait varier l'échelle de Gibbs et maintient fixe l'hamiltonien.
La variation radiale de~\eqref{vii:eq:diferencial-horizonte} modifie
également \(\mathsf H_R\) et conserve ses propres termes géométriques.

\subsection{Section thermométrique et conservation du registre}
\label{viii:subsec:seccion-termometrica-registro}

L'état \(\rho_0\) et l'hamiltonien de réponse permettent de fixer une
coordonnée thermométrique antérieure à tout lecteur numérique candidat.
Pour le rayon fixé, soit \(\mathcal L_\Theta=\mathbb R u_R\) une
droite orientée dont la section positive \(u_R\) correspond à l'état
de référence. Dans les états fidèles de Gibbs du même hamiltonien,
\begin{equation}
 \rho_\beta=\frac{e^{-\beta\mathsf H_R}}
                  {\operatorname{Tr}e^{-\beta\mathsf H_R}},\qquad
 \Theta_R(\rho_\beta)=\frac{1}{\beta\tau_R}u_R,
 \quad \beta>0.
 \label{viii:eq:lector-termometrico}
\end{equation}
L'identité opératorielle \(-\log\rho_\beta=\beta\mathsf H_R+(\log Z_\beta)I\) détermine \(\beta\) : soustraire deux valeurs propres d'énergies différentes donne leur quotient unique. En particulier, \(\Theta_R(\rho_0)=u_R\). La droite \(\mathcal L_E\) est l'espace unidimensionnel des grandeurs énergétiques.

\Needspace{23\baselineskip}
\begin{theorem}[Transducteur positif et critère d'identification]
\label{viii:thm:identificacion-transductor}
Il existe une unique application linéaire positive
\(B_R:\mathcal L_\Theta\to\mathcal L_E\) avec
\(B_R(u_R)=\tau_R\). Elle satisfait
\begin{equation}
 B_R(\Theta_R(\rho_\beta))=\beta^{-1}.
 \label{viii:eq:transductor-gibbs}
\end{equation}
La section chronologique compatible est
\begin{equation}
 \Theta_{{\rm clk},R}
 =\frac{E_P}{\tau_R\log3}u_R
 =\frac{2\pi R}{L\log3}u_R,
 \qquad
 B_R(\Theta_{{\rm clk},R})=
 \frac{\mathcal T_{I/A}q_{\mathcal A}}{\log3}.
 \label{viii:eq:seccion-cronologica-compatible}
\end{equation}
Pour toute autre application linéaire \(B_*\) entre les mêmes droites,
\begin{equation}
 B_*=B_R
 \quad\Longleftrightarrow\quad
 B_*(\Theta_{{\rm clk},R})=\frac{h}{108t_0\log3}.
 \label{viii:eq:criterio-una-seccion}
\end{equation}
\end{theorem}
\begin{proof}
La section \(u_R\) est une base d'une droite, de sorte que prescrire
son image définit une unique application linéaire ; la positivité de \(\tau_R\)
garantit la positivité de l'application. Substituer
\eqref{viii:eq:lector-termometrico} prouve
\eqref{viii:eq:transductor-gibbs}. L'expression de \(\tau_R\) et
\eqref{viii:eq:composicion-area-termica} donnent
\eqref{viii:eq:seccion-cronologica-compatible}. Enfin,
\(\Theta_{{\rm clk},R}\) est non nulle. Deux applications linéaires d'une
droite coïncident si et seulement si elles coïncident sur cette section.
\end{proof}

Le critère compare des lecteurs dans une section obtenue préalablement à partir de l'état, de la réponse énergétique et de la normalisation tritique. L'application à une coordonnée numérique concrète exige d'évaluer ce lecteur dans la section spécifiée. Les bases dimensionnelles sont transportées conjointement avec les applications et avec les coordonnées thermiques qu'elles expriment.

\begin{proposition}[Transport de la réponse avec mémoire]
\label{viii:prop:transporte-termico-memoria}
Soit \(J\) une isométrie qui conserve les étiquettes complètes du
registre, et considérons son image comme espace récepteur. Pour
\(\mathsf H'_R=J\mathsf H_RJ^*\), \(\rho'_0=J\rho_0J^*\)
et \(\sigma'=J\sigma J^*\), les espérances énergétiques,
les entropies et le bilan~\eqref{viii:eq:balance-libre-termico} sont conservés.
On a également \(J\mathsf H_R=\mathsf H'_RJ\).
\end{proposition}
\begin{proof}
Sur l'image, \(J\) est unitaire et transporte le calcul fonctionnel.
La cyclicité de la trace conserve les espérances ; les mêmes valeurs propres
non nulles conservent l'entropie et l'entropie relative. L'identité
d'entrelacement utilise \(J^*J=I\). Substituer ces égalités dans le bilan
prouve sa conservation.
\end{proof}

La même équivalence admet une autre échelle \(\tau'=a\tau_R\), avec \(a>0\), et une autre origine énergétique \(e'\). Si \(\mathsf H'=e'I+aJ\mathsf H_RJ^*\), alors
\[
 \frac{\mathsf H'-e'I}{\tau'}
 =J\frac{\mathsf H_R}{\tau_R}J^*,\qquad
 \frac{e^{-\mathsf H'/\tau'}}
      {\operatorname{Tr}_{\operatorname{im}J}e^{-\mathsf H'/\tau'}}
 =J\rho_0J^*.
\]
Le facteur scalaire \(e^{-e'/\tau'}\) s'annule dans l'état ;
les différences énergétiques sont multipliées par \(a\), et l'entropie
relative reste invariante. Le bilan correspondant a donc
pour coefficient \(\tau'\). L'espérance énergétique complète
conserve l'origine : \(\langle\mathsf H'\rangle
=e'+a\langle\mathsf H_R\rangle\). La comparaison des réalisations
s'exprime ainsi au moyen de leurs exposants centrés et de leurs échelles propres.
Si l'espérance alimente les lectures \(m=E/c^2\) et
\(r_g=GE/c^4\), maintenir \(c,G\) donne
\[
 m'=\frac{e'}{c^2}+am,\qquad
 r'_g=\frac{Ge'}{c^4}+ar_g.
\]
La normalisation statistique et le transport des grandeurs maintiennent donc des fonctions distinctes pour la même origine enregistrée.

La composition relie l'incrément surfacique de Barbero--Immirzi, l'énergie du calendrier et la distribution angulaire du bord. Son transport retient la provenance sectorielle et l'archive : le retour de phase reste accompagné de l'avancée de la mémoire de l'état enrichi.


\subsection{Coordinatisation de Kruskal et échange des extérieurs}
\label{vii:subsec:kruskal-dos-hojas}

Fixons un état exprimé \(x\), son rayon \(R=R(x)>0\) et les sections positives \(c,G\). Dans la coordinatisation qui suit, \(x,R,c,G\) sont des paramètres fixes ; \(t,r\) sont les coordonnées spatio-temporelles. La réalisation extérieure des deux feuillets est spécifiée par
\begin{equation}
 f(r)=1-\frac Rr,\qquad
 \mathrm ds^2=-f(r)c^2\mathrm dt^2+f(r)^{-1}\mathrm dr^2
                 +r^2\mathrm d\Omega_2^2,\qquad r>R.
 \label{vii:eq:metrica-exterior}
\end{equation}
Pour le lecteur de feuillet \(\sigma\in\{+1,-1\}\), nous définissons
\begin{align}
 r_*&=r+R\log\left(\frac rR-1\right),\\
 \mathcal U&=-\sigma\exp\left(-\frac{ct-r_*}{2R}\right),&
 \mathcal V&=\sigma\exp\left(\frac{ct+r_*}{2R}\right).
 \label{vii:eq:carta-kruskal}
\end{align}
Les symboles \(\mathcal U,\mathcal V\) sont des coordonnées et se
distinguent des opérateurs de transport des routes.

\begin{theorem}[Réalisation géométrique et involution de feuillet]
\label{vii:thm:kruskal-hojas}
La coordinatisation~\eqref{vii:eq:carta-kruskal} réalise les deux feuillets comme les deux extérieurs \(\mathcal U\mathcal V<0\) de l'extension de Kruskal. On a
\begin{equation}
 -\mathcal U\mathcal V=\left(\frac rR-1\right)e^{r/R},\qquad
 \mathrm ds^2=-\frac{4R^3}{r}e^{-r/R}\,
                    \mathrm d\mathcal U\,\mathrm d\mathcal V
                    +r^2\mathrm d\Omega_2^2.
 \label{vii:eq:metrica-kruskal}
\end{equation}
L'involution
\begin{equation}
 \mathcal C_{\rm H}:(\sigma,\mathcal U,\mathcal V)
 \longmapsto(-\sigma,-\mathcal U,-\mathcal V)
 \label{vii:eq:involucion-horizonte}
\end{equation}
est isométrique, fixe le rayon aréolaire, échange les extérieurs et préserve
l'ensemble des horizons \(\mathcal U\mathcal V=0\).
La métrique est régulière en \(r=R\) et satisfait les équations du vide
dans le domaine régulier \(r>0\).
\end{theorem}
\begin{proof}
Multiplier les coordonnées donne la première identité. Comme
\(\mathrm dr_*/\mathrm dr=f^{-1}\),
\[
 \mathrm d\mathcal U=\frac{\mathcal U}{2R}
            (-c\,\mathrm dt+f^{-1}\mathrm dr),\qquad
 \mathrm d\mathcal V=\frac{\mathcal V}{2R}
            ( c\,\mathrm dt+f^{-1}\mathrm dr).
\]
Leur produit reproduit la partie radiale et temporelle de la métrique.
La fonction \(z\mapsto(z-1)e^z\) a pour dérivée \(ze^z>0\)
pour \(z>0\), ce qui détermine \(r\) à partir du produit.
À l'extérieur, on retrouve
\(\sigma=\operatorname{sign}\mathcal V\) et
\(t=(R/c)\log(-\mathcal V/\mathcal U)\).
En \(r=R\), le coefficient radial de la métrique étendue est
\(-4R^2/e\), fini et non nul. L'inversion conserve le
produit, le rayon et \(\mathrm d\mathcal U\mathrm d\mathcal V\),
et échange les deux choix de signe.

Pour vérifier le vide dans la coordinatisation extérieure, les composantes mixtes indépendantes du tenseur d'Einstein se réduisent à
\[
 \mathsf G^t{}_t=\mathsf G^r{}_r
 =\frac{rf'+f-1}{r^2},\qquad
 \mathsf G^\theta{}_\theta=\mathsf G^\phi{}_\phi
 =\frac{f'}r+\frac{f''}2.
\]
Avec \(f'=R/r^2\) et \(f''=-2R/r^3\), toutes s'annulent. L'expression régulière~\eqref{vii:eq:metrica-kruskal} prolonge le résultat à travers les horizons dans le domaine \(r>0\).
\end{proof}

Cette réalisation conserve \(E_{\rm H}=c^4R/(2G)\) et \(S_{\rm H}=k_B\pi R^2/\ell_P^2\), qui dépendent du rayon aréolaire. La température \(T_{\rm cos}\) de \eqref{vii:eq:horizonte-cosmologico} appartient à la réalisation cosmologique qui y est déclarée. La coordinatisation extérieure \eqref{vii:eq:metrica-exterior} conserve sa propre normalisation temporelle. L'application géométrique établit l'échange des extérieurs ; la comparaison avec la formation par effondrement, les spectres de rayonnement ou les modes quasi-normaux utilise en outre leurs conditions physiques respectives.

\subsection{Transport d'états, d'observables et de raffinements}
\label{vii:subsec:transporte-informacion-horizonte}

Notons \(F_x\) la coordinatisation des deux extérieurs qui vient d'être construite. Dans chaque coordinatisation angulaire régulière, son jacobien temporel et radial est
\begin{equation}
 \left|\frac{\partial(\mathcal U,\mathcal V)}{\partial(t,r)}\right|
 =\frac{cr}{2R^3}e^{r/R}>0.
 \label{vii:eq:jacobiano-kruskal}
\end{equation}
L'identité résulte des deux différentielles de la démonstration précédente. Les autres registres de l'état se transportent comme des coordonnées de fibre ; la coordinatisation conserve leur identité et leur mémoire.

\begin{proposition}[Équivalence unitaire des lectures de feuillet et d'extérieur]
\label{vii:prop:transporte-unitario-horizonte}
Soient \(\mu\) une mesure de feuillet et \(\nu=(F_x)_*\mu\) sa mesure
transportée. Le changement de variables définit un opérateur unitaire
\begin{equation}
 W_{\rm H}:L^2(\mu)\longrightarrow L^2(\nu),\qquad
 (W_{\rm H}\psi)(z)=\psi(F_x^{-1}z).
 \label{vii:eq:unitaria-horizonte}
\end{equation}
Avec des densités par rapport à des mesures de coordonnées fixées, la même application incorpore le facteur jacobien correspondant. La conjugaison
\begin{equation}
 \mathfrak I_{\rm H}(B)=W_{\rm H}BW_{\rm H}^*
 \label{vii:eq:mapa-observables-horizonte}
\end{equation}
est un \(*\)-isomorphisme unital complètement positif. Il transporte l'
état, ses espérances et son calcul fonctionnel modulaire. Pour la mesure
symétrique des deux feuillets, il entrelace aussi leurs involutions.
\end{proposition}
\begin{proof}
La définition de la mesure image donne \(\int|\psi\circ F_x^{-1}|^2\,\mathrm d\nu
=\int|\psi|^2\,\mathrm d\mu\). La coordinatisation est inversible dans les extérieurs, de sorte que l'isométrie est surjective. La conjugaison unitaire préserve le produit, l'adjoint et l'identité ; dans chaque amplification matricielle, elle conjugue par \(I_n\otimes W_{\rm H}\), ce qui démontre la positivité complète. Pour un état \(\rho\), la cyclicité de la trace donne \(\operatorname{Tr}(W_{\rm H}\rho W_{\rm H}^*
\,\mathfrak I_{\rm H}(B))=\operatorname{Tr}(\rho B)\). Le calcul fonctionnel implique \(g(W_{\rm H}\rho W_{\rm H}^*)=W_{\rm H}g(\rho)W_{\rm H}^*\) dans son domaine, en particulier pour le logarithme et le flot modulaire d'un état fidèle sur son support. L'égalité \(F_x\circ\mathcal C_{\rm hoja}=\mathcal C_{\rm H}\circ F_x\) démontre l'entrelacement des involutions ; la mesure symétrique les réalise unitairement.
\end{proof}

Si les raffinements conservent \((\sigma,t,r,R)\) et utilisent la même
isométrie sur les registres de fibre, les inclusions
\(j_m^{\rm hoja}\) et \(j_m^{\rm ext}\) satisfont
\begin{equation}
 W_{{\rm H},m+1}j_m^{\rm hoja}
 =j_m^{\rm ext}W_{{\rm H},m}.
 \label{vii:eq:naturalidad-horizonte}
\end{equation}
En effet, le changement de variables agit sur les coordonnées géométriques
conservées et l'inclusion sur les mêmes coordonnées de fibre ; évaluer
les deux compositions donne la même fonction et la même densité transportée.
L'égalité étend l'application au système raffiné sans réinitialiser les registres.

Pour une bipartition \(\mathcal H_A\otimes\mathcal H_B\), un
transport factorisé \(U_A\otimes U_B\) satisfait
\[
 \rho'_A=U_A\rho_AU_A^*,\qquad s(\rho'_A)=s(\rho_A).
\]
La première identité découle de la trace partielle et la seconde du spectre.
Pour une transformation unitaire générale, la formulation équivalente transporte également
l'algèbre du sous-système :
\(\mathcal B(\mathcal H_A)\otimes I\mapsto
U(\mathcal B(\mathcal H_A)\otimes I)U^*\).
On conserve ainsi la même question opérationnelle concernant l'information
accessible. Les égalités d'entropie s'appliquent aux états dont les entropies
pertinentes sont finies.

La relation entre information et gravitation s'exprime par des opérations distinctes et composables : l'état et la bipartition fixent l'information accessible ; l'incidence et Barbero--Immirzi fixent la résolution d'aire ; l'action et la gravitation construisent l'unité \(\ell_P^2\) ; Boltzmann réalise dimensionnellement l'entropie ; la coordinatisation d'horizon transporte la géométrie et ses observables. Les lois finies conservent l'entropie relative, et les réalisations géométriques conservent leurs mesures, leur normalisation temporelle et leurs domaines.
